\documentclass[twocolumn]{aastex631}
\graphicspath{{../figures/}}
\begin{document}

\title{Optical denominator for gas-fraction versus efficiency tests: an SDSS DR17 pilot}
\shorttitle{Optical denominator for gas-fraction versus efficiency tests}
\author{NebulaMind Research Autopilot}
\affiliation{Local reproducible pilot run; public SDSS DR17 data only}

\begin{abstract}
We execute a bounded actual-data pilot for the Galaxy Evolution research topic ``Distinguishing molecular-gas depletion from suppressed star-formation efficiency in quenched galaxies''.  The analysis uses a public SDSS DR17 emission-line galaxy sample with stellar-mass, specific-SFR, photometry, and optical emission-line measurements.  The pilot question is: How many massive quenched or transitioning SDSS galaxies with valid emission-line measurements are available as a denominator for CO gas-fraction/depletion-time follow-up?  The result is intended as a survey denominator or proxy measurement, not as a full causal test of feedback physics.
\end{abstract}

\section{Scope}
This manuscript is an AAS-style pilot generated from a research proposal.  It uses actual public survey data and preserves the analysis artifacts, but it deliberately stays inside the observables available in the SDSS sample.  The full proposal requires CO or dust-based molecular gas masses, aperture-matched SFRs, morphology, and environment labels.

\section{Data and Measurements}
The input table is the SDSS DR17 emission-line sample from run SDSS-AGN-SFR-PILOT-20260708T122000Z.  It contains 60,000 galaxies after requiring spectroscopic galaxy classification, redshift 0.02--0.12, finite stellar-mass and specific-SFR estimates, and signal-to-noise at least 3 in H$\alpha$, H$\beta$, [O~III] $\lambda5007$, and [N~II] $\lambda6584$.  BPT classes are recomputed from the line ratios using the Kauffmann et al. and Kewley et al. demarcations.  A local-density ranking is computed from the 10th nearest neighbour in approximate comoving Cartesian coordinates and is used only as an internal density proxy.

\section{Pilot Result}
\begin{itemize}
\item The massive transition/quenched denominator contains 6,729 galaxies in the SDSS emission-line sample.
\item Its optical BPT AGN fraction is 0.549; median log H-alpha luminosity proxy is 40.06.
\item The median H-alpha luminosity proxy is -0.66 dex offset from massive star-forming emission-line galaxies.
\item SDSS optical data cannot distinguish molecular-gas depletion from reduced star-formation efficiency; this paper identifies the CO follow-up denominator and optical baseline.
\end{itemize}

\begin{figure}
\centering
\includegraphics[width=\columnwidth]{m3\_p2\_gas\_depletion\_efficiency\_figure1.pdf}
\caption{Topic-specific SDSS DR17 pilot measurement. The figure is a proxy or denominator diagnostic, not the full multi-survey test.}
\end{figure}

\section{Interpretation Guard}
This pilot should not be read as a completed proof of the full feedback proposal.  It is a reproducible SDSS measurement that defines a denominator, proxy, or validation target for the larger research design.  In particular, it does not by itself establish causal AGN feedback, gas escape, radio-jet coupling efficiency, X-ray maintenance heating, molecular-gas depletion, or simulation correctness.

\section{Reproducibility}
Run identifier: SDSS\_REMAINING\_TOPIC\_PILOTS\_20260708T125828Z.  The run directory preserves the topic-specific JSON summary, figure files, manuscript source, compile log, and compiled PDF.  The workflow used read-only public SDSS-derived data already cached from the first pilot plus local artifact writes only.

\acknowledgments
This pilot used public SDSS DR17 data products and open-source Python tools.

\begin{thebibliography}{}
\bibitem[Baldwin et al.(1981)]{baldwin1981} Baldwin, J.~A., Phillips, M.~M., \& Terlevich, R. 1981, PASP, 93, 5
\bibitem[Kauffmann et al.(2003)]{kauffmann2003} Kauffmann, G., Heckman, T.~M., Tremonti, C., et al. 2003, MNRAS, 346, 1055
\bibitem[Kewley et al.(2001)]{kewley2001} Kewley, L.~J., Dopita, M.~A., Sutherland, R.~S., Heisler, C.~A., \& Trevena, J. 2001, ApJ, 556, 121
\bibitem[York et al.(2000)]{york2000} York, D.~G., Adelman, J., Anderson, J.~E., Jr., et al. 2000, AJ, 120, 1579
\end{thebibliography}

\end{document}
